\originalchapter{连接图、virial与液气相变}
\label{ch:clusters}
\source{8.333 L15--L18及L19的课历回顾；旧卷相关公式在此用于连接图与临界问题。}
\originalsection{Mayer函数与连接展开}
两体势$v_{ij}$使$\e^{-\beta\sum v_{ij}}=\prod_{i<j}(1+f_{ij})$，$f_{ij}=\e^{-\beta v_{ij}}-1$。每项可画为粒子顶点与所选$f$边的图。构型积分中不相连的图分解成乘积，而$\log\Xi$只保留连接图，这是累积量与连接相关的同一组合结构。

写$\beta p=\sum_{l\ge1}b_lz^l$、$n=\sum l b_lz^l$，$z$无量纲。平移不变、短程势下
\begin{align}
b_1&=\lambda_T^{-d},\quad b_2=\frac1{2\lambda_T^{2d}}\int f(r)\dd^dr,\\
b_3&=\frac1{6\lambda_T^{3d}}\iint[ f_{12}f_{13}+f_{12}f_{23}+f_{13}f_{23}+f_{12}f_{13}f_{23}]\dd^dr_2\dd^dr_3,
\end{align}
固定$r_1=0$后总体积已除去。三个树图与一个三角图皆为连接，不能只留三角便称为$b_3$。
\begin{center}\begin{tikzpicture}[scale=.8]
\foreach \x in {0,3.4,6.8}{\fill(\x,0)circle(2pt);\fill(\x+1,1)circle(2pt);\fill(\x+2,0)circle(2pt);\draw(\x,0)--(\x+1,1)--(\x+2,0);}
\draw(6.8,0)--(8.8,0);\node[below]at(1,-.3){连接树};\node[below]at(4.4,-.3){另一树（标签不同）};\node[below]at(7.8,-.3){不可约三角};
\end{tikzpicture}\end{center}
图中的第二树仅示形状；实际标签置换在代数中给三份，不能把形状数当带标签图数。
\originalsection{消去逸度与virial系数}
由数密度级数反解$z=n/b_1-2b_2n^2/b_1^3+[8b_2^2/b_1^5-3b_3/b_1^4]n^3+\cdots$，回代压力得到
\begin{equation}
\beta p=n+B_2n^2+B_3n^3+\cdots,\quad B_2=-\frac{b_2}{b_1^2},\quad B_3=\frac{4b_2^2}{b_1^4}-\frac{2b_3}{b_1^3}.
\end{equation}
因此$B_2=-\frac12\int f$，$B_3=-\frac13\iint f_{12}f_{23}f_{31}$，树图在由逸度换成密度时取消。这个取消解释为何virial展开只含相应二连接图，而不是随意删掉的经验规则。

三维硬球直径$\sigma$有$f=-1$于$r<\sigma$、其余0，$B_2=2\pi\sigma^3/3$，$B_3=5\pi^2\sigma^6/18$。后者为三个两两距离小于$\sigma$的几何体积，需做重叠球积分；可取$r_{12}=r$，两个半径$\sigma$球交体积$\pi(4\sigma+r)(2\sigma-r)^2/12$，再$\int_0^\sigma4\pi r^2\dd r$并除3，得到所给值。这样保留第三virial的可复算主计算。
\originalsection{van der Waals与共存}
排斥硬芯给$b=2\pi\sigma^3/3$，弱吸引外尾把$f\simeq-\beta v$，$B_2=b-a/(\kb T)$，$a=-2\pi\int_\sigma^\infty r^2v(r)\dd r>0$。把排斥近似为有效自由体积、吸引近似为平均密度作用，得到$p=n\kb T/(1-bn)-an^2$。这个重组比仅截断$B_2$更能表现密度非线性，却不是精确簇展开到所有阶。

临界点满足$\partial_np=\partial_n^2p=0$，求得$n_c=1/(3b)$、$\kb T_c=8a/(27b)$、$p_c=a/(27b^2)$。低于$T_c$曲线有$\partial_np<0$的不稳定区。平衡是自由能的凸包，两相须同$T,p,\mu$；对摩尔体积作Maxwell等面积构造正是相同化学势条件的积分表达，而非把负压缩率支保留下来。
\originalsection{平均场指数与临界涨落的入口}
临界附近取密度差$\phi=n-n_c$，消去平移后自由能近$F/V=f_0+r\phi^2/2+u\phi^4/4-h\phi$，$u>0$。极小条件$r\phi+u\phi^3=h$给$\phi\propto(-r)^{1/2}$、$\chi\propto r^{-1}$、临界$\phi\propto h^{1/3}$，即$\beta=1/2,\gamma=1,\delta=3$，热容有有限跳跃$\alpha=0$。这些指数忽略空间涨落；第二编将通过Gaussian关联、Ginzburg准则与RG说明何时可用、何时需要新的固定点。一般临界行为不能仅由局部势能的一阶极小化决定。
